% =====================================================================
%  Chess IR — Guide du participant
%  Règles du championnat et utilisation de l'application chessir.vercel.app
%  Compilation : latexmk -pdf guide-chess-ir.tex
% =====================================================================
\documentclass[11pt,a4paper]{article}

\usepackage[T1]{fontenc}
\usepackage[utf8]{inputenc}
\usepackage[french]{babel}
\usepackage{lmodern}
\usepackage[scaled=0.95]{helvet}
\renewcommand{\familydefault}{\sfdefault}
\usepackage[a4paper,margin=2.2cm,headheight=14pt]{geometry}
\usepackage[table]{xcolor}
\usepackage{graphicx}
\usepackage{tikz}
\usetikzlibrary{positioning}
\usepackage{titlesec}
\usepackage{enumitem}
\usepackage{booktabs}
\usepackage{tabularx}
\usepackage{array}
\usepackage{fancyhdr}
\usepackage{fontawesome5}
\usepackage[most]{tcolorbox}
\usepackage[hidelinks]{hyperref}

% ---------------------------------------------------------------------
% Identité visuelle Chess IR
% ---------------------------------------------------------------------
\definecolor{nuit}{HTML}{0C0F2C}
\definecolor{nuitclair}{HTML}{1A1E4A}
\definecolor{or}{HTML}{E09A24}
\definecolor{orclair}{HTML}{F5B83D}
\definecolor{violet}{HTML}{6C5CE7}
\definecolor{vert}{HTML}{16A34A}
\definecolor{rouge}{HTML}{DC2626}
\definecolor{gris}{HTML}{5B6185}
\definecolor{fond}{HTML}{F6F6FB}

\hypersetup{colorlinks=true,linkcolor=nuit,urlcolor=or}

\titleformat{\section}{\Large\bfseries\color{nuit}}{\colorbox{orclair}{\makebox[1.6em]{\color{nuit}\thesection}}}{0.7em}{}
\titleformat{\subsection}{\large\bfseries\color{nuitclair}}{\thesubsection}{0.6em}{}
\titlespacing*{\section}{0pt}{1.6em}{0.8em}

\setlist{itemsep=0.25em,topsep=0.4em}
\setlist[itemize]{label=\textcolor{or}{\small\faChessPawn}}

\pagestyle{fancy}
\fancyhf{}
\fancyhead[L]{\small\color{gris}Chess IR — Guide du participant}
\fancyhead[R]{\small\color{gris}\href{https://chessir.vercel.app}{chessir.vercel.app}}
\fancyfoot[C]{\small\color{gris}\thepage}
\renewcommand{\headrulewidth}{0.4pt}
\renewcommand{\headrule}{\hbox to\headwidth{\color{orclair}\leaders\hrule height \headrulewidth\hfill}}

% Encadrés
\newtcolorbox{aretenir}[1][À retenir]{enhanced,colback=fond,colframe=or,boxrule=0.8pt,arc=3mm,
  fonttitle=\bfseries,coltitle=nuit,colbacktitle=orclair,title={\faStar\ #1},left=3mm,right=3mm}
\newtcolorbox{attention}[1][Attention]{enhanced,colback=red!3,colframe=rouge,boxrule=0.8pt,arc=3mm,
  fonttitle=\bfseries,colbacktitle=rouge,title={\faExclamationTriangle\ #1},left=3mm,right=3mm}
\newtcolorbox{astuce}[1][Astuce]{enhanced,colback=violet!4,colframe=violet,boxrule=0.8pt,arc=3mm,
  fonttitle=\bfseries,colbacktitle=violet,title={\faLightbulb\ #1},left=3mm,right=3mm}

% Étapes numérotées
\newcounter{etape}
\newcommand{\etape}[1]{\stepcounter{etape}\par\noindent
  \tikz[baseline=-0.6ex]\node[circle,fill=orclair,text=nuit,font=\bfseries\small,inner sep=2.2pt,minimum size=1.5em]{\theetape};\ \textbf{#1}\par\smallskip}
\newcommand{\resetetapes}{\setcounter{etape}{0}}

\newcommand{\app}{\textbf{Chess~IR}}
\newcommand{\bouton}[1]{\tikz[baseline=(b.base)]\node[draw=or,fill=orclair!25,rounded corners=2pt,inner sep=2pt,font=\small\bfseries](b){#1};}
\newcommand{\lien}{\href{https://chessir.vercel.app}{\textbf{chessir.vercel.app}}}

\begin{document}

% =====================================================================
% Page de titre
% =====================================================================
\begin{titlepage}
\begin{tikzpicture}[remember picture,overlay]
  \fill[nuit] (current page.south west) rectangle (current page.north east);
  \shade[inner color=violet!55!nuit,outer color=nuit] ([yshift=-1cm]current page.north) circle (11cm);
  % Damier en perspective
  \begin{scope}[shift={([yshift=4.2cm]current page.south)}]
    \foreach \i in {0,...,5}{
      \foreach \j in {-7,...,6}{
        \pgfmathtruncatemacro{\parite}{mod(\i+\j+20,2)}
        \ifnum\parite=0
          \pgfmathsetmacro{\ya}{-\i*0.75}\pgfmathsetmacro{\yb}{-(\i+1)*0.75}
          \pgfmathsetmacro{\sa}{0.9+\i*0.32}\pgfmathsetmacro{\sb}{0.9+(\i+1)*0.32}
          \fill[white,opacity=0.06] (\j*\sa,\ya) -- ({(\j+1)*\sa},\ya) -- ({(\j+1)*\sb},\yb) -- (\j*\sb,\yb) -- cycle;
        \fi
      }
    }
  \end{scope}
  \node at ([yshift=-8.2cm]current page.north) {\includegraphics[width=5cm]{../icone.png}};
  \node[text=white,font=\fontsize{46}{50}\selectfont\bfseries] at ([yshift=-12.6cm]current page.north)
    {Chess \textcolor{orclair}{IR}};
  \node[text=white!75,font=\Large] at ([yshift=-14cm]current page.north) {Jouer. Progresser. Grimper.};
  \node[text=white,font=\LARGE\bfseries] at ([yshift=-17cm]current page.north) {Guide du participant};
  \node[text=white!80,font=\large,align=center] at ([yshift=-18.4cm]current page.north)
    {Règles du championnat d'échecs de la classe\\et utilisation de l'application};
  \node[text=orclair,font=\Large\bfseries] at ([yshift=-21.2cm]current page.north)
    {\faGlobe\ \href{https://chessir.vercel.app}{chessir.vercel.app}};
\end{tikzpicture}
\end{titlepage}

\tableofcontents
\newpage

% =====================================================================
\section{Chess IR en bref}
% =====================================================================
\app{} est le championnat d'échecs amical de la classe. Débutants, joueurs occasionnels ou confirmés :
tout le monde joue, personne n'est éliminé, et chacun progresse à son rythme.

\begin{itemize}
  \item \textbf{On joue sur Chess.com}, la plateforme d'échecs en ligne.
  \item \textbf{On gère tout sur \lien} : adversaires de la semaine, déclaration des résultats,
        classement, statistiques, récompenses.
  \item \textbf{On suit les matchs en direct} dans l'\textbf{Arène} de l'application, avec des notifications sur le téléphone.
  \item \textbf{On partage dans le groupe WhatsApp} de la classe : annonces de matchs, matchs en direct, résultats,
        classements et récompenses.
\end{itemize}

\begin{aretenir}
Chaque semaine : \textbf{3 adversaires}, connus dès le dimanche. Tu joues tes parties sur Chess.com quand
ça t'arrange dans la semaine, tu déclares le résultat dans \app{}, ton adversaire le confirme, et tes points
s'ajoutent au classement général. Victoire = 3 points, nul = 1 point, défaite = 0 point.
\end{aretenir}

% =====================================================================
\section{Ce qu'il te faut}
% =====================================================================
\begin{enumerate}[label=\textcolor{or}{\bfseries\arabic*.}]
  \item \textbf{Un compte Chess.com} — \emph{obligatoire}. C'est là que se jouent toutes les parties.
        Il est gratuit : \href{https://www.chess.com}{www.chess.com}. Note bien ton \emph{nom d'utilisateur}
        (pseudo) : il te sera demandé.
  \item \textbf{Un compte \app{}}, créé sur \lien{} (voir section~\ref{sec:inscription}).
  \item \textbf{Ton smartphone}, avec Chrome (Android) ou Safari (iPhone). \app{} est conçu pour le téléphone :
        installe-le sur ton écran d'accueil (section~\ref{sec:installer}) et utilise-le comme une application.
  \item \textbf{Les notifications activées} sur ton téléphone (section~\ref{sec:notifs}) : c'est ainsi que tu sais
        qu'on te défie ou qu'un match passe en direct.
  \item \textbf{L'accès au groupe WhatsApp} de la classe. Le lien du groupe est disponible dans l'application
        (bouton \bouton{Ouvrir le groupe} dans les fenêtres de partage).
\end{enumerate}

\begin{attention}[Ton pseudo Chess.com]
Le pseudo saisi dans \app{} doit être \textbf{exactement} celui de ton compte Chess.com. L'application le vérifie
auprès de Chess.com : un pseudo inexistant est refusé. Une fois le championnat lancé, \textbf{seul
l'administrateur} peut le modifier. C'est grâce à lui que l'application retrouve automatiquement tes parties.
\end{attention}

% =====================================================================
\section{S'inscrire et installer l'application}\label{sec:inscription}
% =====================================================================
\subsection{Créer ton compte}
\resetetapes
\etape{Ouvre \lien{} et choisis l'onglet \bouton{Inscription}.}
Renseigne ton nom ou pseudonyme (celui que verront les autres), ton pseudo Chess.com, ton adresse e-mail et
un mot de passe (6 caractères minimum).

\etape{Confirme ton adresse e-mail si on te le demande.}
Si un e-mail de confirmation t'est envoyé, clique sur le lien qu'il contient, puis connecte-toi.

\etape{Attends ton inscription au championnat.}
Créer un compte ne suffit pas : c'est l'\textbf{administrateur} qui inscrit chaque joueur au championnat.
Tant que ce n'est pas fait, l'accueil affiche « Pas encore inscrit ».

\etape{Vérifie ton profil.}
Dans \bouton{Profil}, ton pseudo Chess.com doit apparaître avec la mention \textcolor{vert}{\faCheckCircle\ vérifié}.
Ton Elo et ton avatar sont repris automatiquement de Chess.com.

\subsection{Installer l'application sur ton téléphone}\label{sec:installer}
\app{} s'installe comme une vraie application, sans passer par un store. C'est la façon recommandée de l'utiliser :
une icône sur ton écran d'accueil, en plein écran, ouverte en un geste.
\begin{itemize}
  \item \textbf{Android (Chrome)} : touche \bouton{Installer l'application} sur l'accueil ou dans ton profil.
  \item \textbf{iPhone (Safari)} : touche \faShareSquare\ \textbf{Partager}, puis \textbf{Sur l'écran d'accueil}, puis \textbf{Ajouter}.
\end{itemize}
Ouvre bien le lien dans Chrome ou Safari, pas dans le navigateur intégré de WhatsApp ou d'Instagram.
L'application installée \textbf{se met à jour toute seule} : rien à télécharger de nouveau.

\subsection{Activer les notifications}\label{sec:notifs}
Les notifications arrivent sur ton téléphone \textbf{même quand l'application est fermée} : défi reçu, match en direct,
résultat à confirmer, tirage publié, récompense\dots
\resetetapes
\etape{Sur l'accueil, touche \bouton{Activer les notifications}, puis \textbf{Autoriser}.}
\etape{C'est tout.} Tu peux les couper ou les réactiver à tout moment dans \bouton{Profil} $\rightarrow$
\textbf{Notifications sur ce téléphone}.
\begin{attention}[Une activation par téléphone]
\begin{itemize}[leftmargin=1.2em]
  \item S'inscrire ne suffit pas : \textbf{chacun active les notifications lui-même}, une fois, sur son téléphone.
        À refaire si tu changes de téléphone.
  \item \textbf{iPhone} : les notifications ne fonctionnent que depuis l'application \textbf{installée} sur l'écran
        d'accueil (section~\ref{sec:installer}), ouverte depuis son icône.
  \item Si tu as refusé par erreur, l'application t'explique comment les réactiver dans les réglages du téléphone.
\end{itemize}
\end{attention}
\textbf{Son des alertes} : quand l'application est ouverte, un défi ou un match en direct est aussi signalé par un son
et une vibration. Le son est activé par défaut et se coupe dans \bouton{Profil} $\rightarrow$ \textbf{Son des alertes}.

\subsection{Mot de passe oublié}
Sur l'écran de connexion, touche \textbf{Mot de passe oublié ?}, saisis \textbf{l'adresse e-mail de ton compte}
(celle utilisée à l'inscription) et ouvre le lien reçu par e-mail sur ton téléphone : tu choisis ton nouveau mot de passe
et tu es connecté. Le lien ne sert qu'une fois ; pense à regarder dans les spams. Pour une autre adresse que celle de
ton compte, aucun e-mail ne part.

% =====================================================================
\section{Le format du championnat}
% =====================================================================
\begin{itemize}
  \item \textbf{Un championnat dure plusieurs semaines}, autant que l'administrateur le décide, jusqu'à sa clôture finale.
  \item \textbf{Aucune élimination} : tout le monde joue jusqu'au bout, quel que soit son résultat.
  \item \textbf{Classement cumulatif} : les points de chaque semaine s'additionnent ; le classement ne repart jamais à zéro.
  \item \textbf{Système d'appariement de type suisse} : petit à petit, tu affrontes des joueurs qui ont des résultats proches
        des tiens, ce qui rend les matchs équilibrés pour les débutants comme pour les confirmés.
\end{itemize}

\subsection{La semaine type}
\begin{center}
\renewcommand{\arraystretch}{1.35}
\begin{tabularx}{0.92\linewidth}{>{\bfseries\color{nuit}}l X}
\toprule
\rowcolor{orclair!25}Dimanche & Clôture de la semaine précédente (classement figé, récompenses) puis publication du nouveau tirage. Tu reçois une notification. \\
Lundi    & Ronde 1 \\
Mercredi & Ronde 2 \\
Vendredi & Ronde 3 \\
\bottomrule
\end{tabularx}
\end{center}

\begin{aretenir}[Souplesse]
Les jours (lundi, mercredi, vendredi) sont \textbf{indicatifs}. Tu peux jouer chaque partie \textbf{n'importe quand,
dès la publication du tirage} (même le dimanche soir), dans l'ordre que tu veux, en t'arrangeant avec ton adversaire selon vos cours. Le seul impératif : que les trois résultats
soient déclarés et confirmés \textbf{avant la clôture du dimanche}.
\end{aretenir}

\subsection{Comment sont choisis tes adversaires}
\begin{itemize}
  \item \textbf{Semaine 1} : tirage au sort.
  \item \textbf{Ensuite} : les joueurs sont rapprochés selon leurs points au classement général.
  \item Tu ne rencontres jamais deux fois le même adversaire dans une semaine, et l'application évite autant que possible
        les adversaires déjà affrontés les semaines précédentes.
  \item Les \textbf{couleurs} (Blancs / Noirs) sont équilibrées au fil des semaines et imposées par le tirage.
  \item \textbf{Nombre impair de joueurs} : à chaque ronde, un joueur est \textbf{exempt} (il ne joue pas). Ce n'est jamais
        le même deux fois dans la semaine, et c'est en priorité celui qui a été le moins souvent exempté.
        Les points accordés à un exempt (0, 1 ou 3) sont fixés par l'administrateur avant le lancement et rappelés sous le classement.
\end{itemize}
Une fois publié, le tirage est \textbf{verrouillé} : il ne change plus pendant la semaine.

% =====================================================================
\section{Jouer un match}
% =====================================================================
\resetetapes
\etape{Ouvre ton match dans \app{}.}
Sur l'accueil, la carte \textbf{« Mon prochain match »} t'indique ton adversaire, ta couleur et le jour prévu.

\etape{Défie ton adversaire depuis \app{}.}
Touche \bouton{Défier} : Chess.com s'ouvre sur ton adversaire pour lancer le défi, et ton adversaire reçoit aussitôt
une notification « \emph{Untel vous défie !} ».

\etape{Ton adversaire répond dans les 5 minutes.}
Il touche \bouton{J'y vais} (avec un compte à rebours) : Chess.com s'ouvre chez lui pour accepter ton défi.
Si ton défi Chess.com a déjà expiré, c'est lui qui t'en envoie un nouveau : tu es forcément connecté, puisque tu viens de défier.

\etape{Le match passe « en direct ».}
Dès que vous avez appuyé tous les deux, toute la classe est prévenue et peut suivre le match dans l'\textbf{Arène}
(section~\ref{sec:arene}). Jouez votre partie sur Chess.com.

\begin{astuce}[Défi sans réponse]
Au bout de 5 minutes sans réponse, le défi disparaît : relance-le plus tard. Un défi lancé par erreur, ou un direct qui
n'a pas eu lieu, s'annule avec \bouton{Annuler} ou \bouton{Fausse alerte}. Tu peux aussi annoncer ton match à l'avance
avec \bouton{Annoncer au groupe}.
\end{astuce}

\begin{attention}[Les conditions d'un match de championnat]
Une partie Chess.com compte pour le championnat seulement si elle réunit \textbf{toutes} ces conditions :
\begin{itemize}[leftmargin=1.2em]
  \item \textbf{les deux joueurs du tirage} ;
  \item \textbf{la cadence du championnat} (par défaut \textbf{Rapide}, par exemple 10 ou 15 minutes) ;
  \item une partie \textbf{classée} (« rated ») ;
  \item \textbf{les couleurs du tirage} : au moment du défi sur Chess.com, \textbf{choisis toi-même ta couleur}
        (Blancs ou Noirs) au lieu de « aléatoire ». L'application te le rappelle sous le bouton \bouton{Défier} ;
  \item échecs \textbf{classiques} (pas de variantes comme le Chess960), jouée entre la publication du tirage et la clôture
        de la semaine.
\end{itemize}
\end{attention}

\begin{aretenir}[Parties amicales]
Toute autre partie entre vous est une \textbf{partie amicale} : l'application l'ignore. Pour un avant-match ou une revanche
pour le plaisir, jouez \textbf{en blitz}, \textbf{en non classé} ou \textbf{couleurs inversées} : elle ne risque pas
d'être prise pour votre match officiel.
\end{aretenir}

% =====================================================================
\section{Déclarer et valider le résultat}
% =====================================================================
Un résultat ne compte que lorsqu'il est \textbf{validé par les deux joueurs} (ou, en cas de litige, par l'administrateur).

\subsection{Avec la détection automatique (le plus simple)}
Après la partie, rouvre le match dans \app{} : l'application recherche votre partie sur Chess.com.
\begin{itemize}
  \item Elle te propose \textbf{la première partie qui remplit toutes les conditions} (« Match de championnat trouvé :
        Victoire ») avec le lien de la partie. Touche \bouton{Accepter ce résultat}.
  \item Si vous avez joué plusieurs parties conformes (une revanche, par exemple), les autres restent proposées :
        touche \bouton{C'est plutôt celle-ci} sur la bonne.
  \item Ton adversaire reçoit une notification et \textbf{confirme} à son tour. Le résultat devient officiel.
  \item Si rien n'apparaît encore, touche \bouton{Rechercher} un peu plus tard (une recherche toutes les 2 minutes au plus) :
        Chess.com peut publier les parties avec du retard, parfois plusieurs heures.
\end{itemize}

\subsection{Déclaration manuelle (toujours possible)}
Inutile d'attendre la détection. Dans le match, choisis \bouton{Victoire}, \bouton{Match nul} ou \bouton{Défaite}
(de ton point de vue), puis \bouton{Confirmer}. Ton adversaire doit ensuite confirmer.

\subsection{Le circuit de validation}
\begin{center}
\begin{tikzpicture}[node distance=0.7cm and 1.25cm,
  st/.style={draw=or,fill=orclair!20,rounded corners=4pt,align=center,font=\small,minimum height=1.05cm,text width=2.35cm},
  ok/.style={st,draw=vert,fill=vert!12}, ko/.style={st,draw=rouge,fill=rouge!8},
  >=stealth,thick]
  \node[st] (a) {\textbf{À jouer}};
  \node[st,right=of a] (b) {\textbf{En attente}\\{\scriptsize résultat déclaré}};
  \node[ok,right=of b] (c) {\textbf{Validé}\\{\scriptsize points attribués}};
  \node[ko,below=of b] (d) {\textbf{Contesté}\\{\scriptsize l'admin arbitre}};
  \draw[->,color=nuit] (a) -- node[above,font=\scriptsize]{déclaration} (b);
  \draw[->,color=nuit] (b) -- node[above,font=\scriptsize]{confirmation} (c);
  \draw[->,color=rouge] (b) -- node[right,font=\scriptsize]{contestation} (d);
  \draw[->,color=nuit] (d.east) -| node[below,pos=0.25,font=\scriptsize]{décision} (c.south);
\end{tikzpicture}
\end{center}

\begin{itemize}
  \item \textbf{Tu as déclaré par erreur ?} Tant que ton adversaire n'a pas répondu, touche \bouton{Retirer ma déclaration}.
  \item \textbf{Tu n'es pas d'accord avec le résultat déclaré par ton adversaire ?} Touche \bouton{Contester} et explique
        pourquoi. Le résultat reste en attente jusqu'à la décision de l'administrateur.
  \item Si la déclaration de ton adversaire ne correspond pas à la partie trouvée sur Chess.com, l'application t'avertit.
\end{itemize}

\begin{attention}[Avant la clôture du dimanche]
Un match qui n'a été ni déclaré ni confirmé à la clôture est enregistré \textbf{« non joué »} : \textbf{0 point pour chacun},
sans vainqueur. Une semaine ne peut pas être clôturée tant qu'un résultat est en attente ou contesté : pense à confirmer
rapidement les résultats de tes adversaires.
\end{attention}

% =====================================================================
\section{Points et classement}
% =====================================================================
\begin{center}
\renewcommand{\arraystretch}{1.35}
\begin{tabular}{l c}
\toprule
\textbf{Résultat} & \textbf{Points} \\
\midrule
\textcolor{vert}{\faTrophy}\ Victoire & \textbf{3} \\
\textcolor{violet}{\faHandshake}\ Match nul & \textbf{1} \\
\textcolor{rouge}{\faTimes}\ Défaite & \textbf{0} \\
Match non joué & 0 (pour chacun) \\
Exempt & fixé par l'administrateur (0, 1 ou 3) \\
\bottomrule
\end{tabular}
\end{center}
\textit{Exemple :} deux victoires et un nul dans la semaine $= 3 + 3 + 1 = 7$ points.

\subsection{En cas d'égalité de points}
Les joueurs sont départagés dans cet ordre :
\begin{enumerate}[label=\textcolor{or}{\bfseries\arabic*.}]
  \item le \textbf{nombre de victoires} ;
  \item le \textbf{Buchholz} : la somme des points des adversaires que tu as affrontés. Battre des adversaires solides compte davantage.
\end{enumerate}
Si l'égalité persiste, les joueurs partagent le même rang.

\subsection{Les deux classements}
\begin{itemize}
  \item \textbf{Général} : cumul de toutes les semaines. Les flèches \textcolor{vert}{\faArrowUp}\,/\,\textcolor{rouge}{\faArrowDown}
        indiquent les places gagnées ou perdues depuis la semaine précédente.
  \item \textbf{Semaine} : les points marqués pendant une semaine donnée (utile pour le titre de Champion de la semaine).
\end{itemize}
\begin{aretenir}
Le classement ne tient compte \textbf{que des résultats validés} dans \app{}. Ton Elo et tes statistiques Chess.com
\textbf{n'influencent jamais} tes points ni ton classement.
\end{aretenir}

% =====================================================================
\section{Récompenses}
% =====================================================================
Les récompenses valorisent la participation et la progression, pas seulement les premiers du classement.

\subsection{Distinctions de la semaine}
Attribuées automatiquement \textbf{à la clôture de chaque semaine} (plusieurs gagnants possibles en cas d'égalité).
\begin{center}
\renewcommand{\arraystretch}{1.3}
\begin{tabularx}{\linewidth}{>{\bfseries}l X}
\toprule
\faCrown\ Champion de la semaine & Le meilleur score de la semaine sur les trois rondes. \\
\faBolt\ Chevalier surprise & La victoire la plus inattendue : contre un joueur classé au moins 3 places au-dessus
  au moment du tirage (en semaine 1 : contre un joueur ayant au moins 200 points d'Elo de plus). \\
\faChartLine\ Meilleure progression & La plus forte remontée au classement par rapport à la semaine précédente. \\
\faFire\ Série gagnante & La plus longue série de victoires en cours (3 minimum). \\
\bottomrule
\end{tabularx}
\end{center}

\subsection{Badges}
Débloqués \textbf{immédiatement}, dès qu'un résultat est validé, et gardés pour toujours.
\begin{center}
\renewcommand{\arraystretch}{1.3}
\begin{tabularx}{\linewidth}{>{\bfseries}l X}
\toprule
Premier pas & Disputer sa première partie officielle. \\
Première victoire & Remporter sa première partie. \\
Semaine complète & Jouer ses trois matchs dans la même semaine. \\
Cinq victoires & Atteindre cinq victoires au total. \\
Trois d'affilée & Enchaîner trois victoires consécutives. \\
Tombeur du leader & Battre le joueur classé premier (à partir de la semaine 2). \\
Première remontée & Gagner des places au classement d'une semaine à l'autre. \\
\bottomrule
\end{tabularx}
\end{center}
Toutes tes récompenses sont visibles dans \bouton{Profil} $\rightarrow$ \textbf{Distinctions}, et celles de tout le monde dans
la page \textbf{Récompenses}.

% =====================================================================
\section{L'application, écran par écran}
% =====================================================================
La barre en bas de l'écran donne accès aux cinq espaces principaux ; l'\textbf{Arène} est le bouton doré au centre.
\begin{center}
\renewcommand{\arraystretch}{1.4}
\begin{tabularx}{\linewidth}{>{\raggedright\arraybackslash\bfseries\color{nuit}}p{3.1cm} X}
\toprule
\faHome\ Accueil & La semaine en cours, ta place et tes points, ton évolution, et la carte \textbf{« Mon prochain match »}
  avec les boutons \bouton{Défier} et \bouton{Déclarer}. Une bannière te signale les résultats à confirmer. \\
\faChessKnight\ Matchs & \textbf{Semaine} : tes trois matchs, leur couleur, leur jour et leur statut (À jouer, En attente,
  Validé, Contesté, Non joué). \textbf{Historique} : tous tes résultats. Une pastille dorée sur l'onglet indique qu'une
  action t'attend. \\
\faFire\ Arène & Les matchs \textbf{en direct} de toute la classe, les défis en cours, et les matchs terminés dans la
  journée (section~\ref{sec:arene}). Le bouton pulse avec un compteur rouge quand il se passe quelque chose. \\
\faTrophy\ Classement & Classement général et classement de chaque semaine. Touche un joueur pour voir son profil. \\
\faUser\ Profil & Tes statistiques du championnat, ta progression, tes distinctions, tes derniers matchs, et un encadré
  séparé \textbf{« Sur Chess.com — hors championnat »} avec ton Elo et ton activité globale sur Chess.com. \\
\faBell\ Notifications & Défi reçu, match en direct, nouveau tirage, résultat à confirmer, résultat confirmé ou contesté,
  décision de l'administrateur, partie trouvée sur Chess.com, semaine clôturée, récompense obtenue. Elles arrivent aussi sur
  ton téléphone si tu as activé les notifications (section~\ref{sec:notifs}). \\
\bottomrule
\end{tabularx}
\end{center}

\subsection{L'Arène : les matchs en direct}\label{sec:arene}
\begin{itemize}
  \item \textbf{Défi reçu} : en haut de l'Arène, avec le compte à rebours et le bouton \bouton{J'y vais}. Si un défi
        t'attend quand tu ouvres l'application, elle s'ouvre directement sur l'Arène.
  \item \textbf{En direct} : une carte face-à-face par match (joueurs, rang, Elo, couleurs, durée). Touche
        \bouton{Regarder} sous un joueur pour ouvrir son profil Chess.com et suivre la partie en cours.
  \item \bouton{Partager le direct au groupe} : tout le monde peut inviter la classe à regarder (section suivante).
  \item \textbf{Défis lancés} : les défis en attente de réponse entre d'autres joueurs.
  \item \textbf{Terminés aujourd'hui} : le vainqueur de chaque direct, avec sa couronne.
  \item Sans match en direct : les rencontres qui restent à jouer cette semaine et les derniers résultats.
\end{itemize}
Un match en direct s'arrête tout seul dès que le résultat est déclaré ou retrouvé sur Chess.com (au plus tard au bout de 90 minutes).

\subsection{Partager dans le groupe WhatsApp}
Comme sur Duolingo, l'application prépare pour toi une \textbf{image} et un \textbf{message} prêts à envoyer :
\begin{itemize}
  \item l'annonce d'un match ou son résultat (depuis le match) ;
  \item un \textbf{match en direct}, avec les liens pour le regarder sur Chess.com (depuis l'Arène, accessible à tous) ;
  \item le classement général ou de la semaine (icône \faShare*\ en haut du classement) ;
  \item une distinction ou un badge (page Récompenses) ;
  \item le tirage de la semaine (par l'administrateur).
\end{itemize}
Touche \bouton{Partager}, choisis \textbf{WhatsApp}, puis le \textbf{groupe de la classe}. Le bouton
\bouton{Ouvrir le groupe} ouvre directement le groupe.

\begin{astuce}
Sur certains téléphones, WhatsApp n'envoie que l'image, sans le texte. Touche alors \bouton{Copier le message}
et colle-le sous l'image avant d'envoyer.
\end{astuce}

% =====================================================================
\section{Chess.com et Chess IR}
% =====================================================================
\app{} lit uniquement les informations \textbf{publiques} de Chess.com ; il ne modifie jamais rien sur ton compte
et ne joue pas à ta place.
\begin{itemize}
  \item \textbf{Ton profil} : avatar, Elo par cadence et bilan global sont synchronisés au plus une fois par jour
        (bouton \bouton{Actualiser} dans ton profil).
  \item \textbf{Ton Elo} affiché dans \app{} est ton Elo Chess.com dans la cadence du championnat (Rapide par défaut).
        Ton \textbf{Elo initial} est celui enregistré lors de ton inscription au championnat.
  \item \textbf{Tes parties} : l'application retrouve la partie de chaque match pour te proposer le résultat et
        te permettre de la revoir sur Chess.com depuis ton historique.
  \item Les données Chess.com sont présentées à part et \textbf{ne comptent jamais} dans le championnat.
\end{itemize}

% =====================================================================
\section{Situations particulières}\label{sec:cas}
% =====================================================================
\begin{center}
\renewcommand{\arraystretch}{1.4}
\begin{tabularx}{\linewidth}{>{\raggedright\arraybackslash\bfseries}p{4.2cm} X}
\toprule
Mon adversaire ne répond pas & Contacte-le dans le groupe. Si le match ne peut pas avoir lieu, préviens l'administrateur :
  il peut le déclarer non joué ou prendre une autre décision. Il n'y a \textbf{aucune victoire automatique}. \\
Déconnexion pendant la partie & Chess.com peut l'enregistrer comme abandon ou perte au temps. Le résultat te sera proposé
  tel quel, mais rien n'est validé sans confirmation : en cas de désaccord, contestez et l'administrateur tranche. \\
Plusieurs parties entre nous & Seules les parties conformes sont proposées : la première est suggérée, les autres restent
  au choix (\bouton{C'est plutôt celle-ci}). Rien n'est validé sans votre confirmation à tous les deux. \\
Couleurs inversées, blitz, non classée & C'est une partie amicale : elle n'est pas proposée. Rejouez dans les bonnes
  conditions, ou, en cas de souci, contactez l'administrateur. \\
Mon adversaire ne répond pas au défi & Le défi expire après 5 minutes. Relance-le plus tard ou convenez d'un moment
  dans le groupe. \\
Je ne reçois pas les notifications & Vérifie dans \bouton{Profil} qu'elles sont activées sur ce téléphone ; sur iPhone,
  ouvre l'application depuis son icône sur l'écran d'accueil. \\
Mot de passe oublié & Écran de connexion $\rightarrow$ \textbf{Mot de passe oublié ?}, avec l'adresse e-mail de ton compte. \\
Partie introuvable & Chess.com peut avoir du retard : réessaie plus tard ou déclare le résultat à la main. \\
Semaine déjà clôturée & Plus aucune déclaration n'est possible : contacte l'administrateur. \\
Changer de pseudo Chess.com & Pendant un championnat en cours, seul l'administrateur peut le faire. \\
\bottomrule
\end{tabularx}
\end{center}

% =====================================================================
\section{Fair-play}
% =====================================================================
\begin{itemize}
  \item \textbf{Aucune aide extérieure} pendant une partie : pas de moteur d'échecs, pas d'analyse, pas de conseils.
  \item Déclare toujours le \textbf{résultat réel}, même quand tu perds.
  \item Respecte ton adversaire, dans la partie comme dans le groupe.
  \item En cas de litige, la décision de l'\textbf{administrateur} est finale. Toutes les déclarations, confirmations,
        contestations et décisions sont conservées dans l'historique du match.
\end{itemize}

% =====================================================================
\section{Le rôle de l'administrateur}
% =====================================================================
L'administrateur organise le championnat. Il :
\begin{itemize}
  \item inscrit les joueurs, et peut suspendre une participation (le joueur n'est alors plus tiré au sort mais reste au classement) ;
  \item fixe les règles avant le lancement (points d'exempt, cadence, fuseau horaire) ; elles sont ensuite figées ;
  \item publie le tirage chaque dimanche et clôture la semaine ;
  \item arbitre les contestations et corrige exceptionnellement un résultat, toujours avec un motif communiqué aux joueurs ;
  \item gère le lien du groupe WhatsApp.
\end{itemize}

% =====================================================================
\section{Questions fréquentes}
% =====================================================================
\begin{description}[style=nextline,leftmargin=0pt,font=\color{nuit}\bfseries]
  \item[Je n'ai pas de compte Chess.com, puis-je participer ?]
    Non : toutes les parties se jouent sur Chess.com. Crée un compte gratuit avant de t'inscrire sur \app{}.
  \item[Je n'ai pas d'Elo sur Chess.com, est-ce grave ?]
    Pas du tout. L'Elo ne sert qu'à l'affichage et à la distinction Chevalier surprise en semaine 1 ; il ne change rien aux points.
  \item[Je suis débutant, vais-je affronter les meilleurs ?]
    Seulement de temps en temps : dès la semaine 2, tu affrontes surtout des joueurs ayant des résultats proches des tiens.
  \item[Puis-je jouer mes trois matchs le même jour ?]
    Oui, si ton adversaire est d'accord, et même dès la publication du tirage. Les jours indiqués sont des repères.
  \item[Une partie pour le plaisir contre mon adversaire peut-elle compter par erreur ?]
    Non si tu la joues en blitz, en non classé ou couleurs inversées : elle est alors ignorée. Et aucune partie ne compte
    sans que vous la confirmiez tous les deux.
  \item[Puis-je regarder les matchs des autres ?]
    Oui : dans l'Arène, touche \bouton{Regarder} sous un joueur en direct. Active les notifications pour être prévenu.
  \item[Mon adversaire a déclaré un mauvais résultat.]
    Touche \bouton{Contester} et explique la situation : l'administrateur verra aussi la partie Chess.com si elle a été trouvée.
  \item[Quand mes points apparaissent-ils ?]
    Dès que le résultat est confirmé par ton adversaire (ou fixé par l'administrateur).
  \item[L'application a changé d'apparence.]
    Elle se met à jour automatiquement après chaque amélioration ; il n'y a rien à faire.
\end{description}

\vfill
\begin{tcolorbox}[enhanced,colback=nuit,colframe=nuit,arc=3mm,halign=center,fontupper=\color{white}]
{\Large\bfseries Chess \textcolor{orclair}{IR}}\\[2pt]
Jouer. Progresser. Grimper.\\[4pt]
\href{https://chessir.vercel.app}{\textcolor{orclair}{\bfseries chessir.vercel.app}}
\end{tcolorbox}

\end{document}
